\begin{problem}{\bf Abundance of primes} (5 points)\\
Suppose that 1024-bit primes are generated by selecting 1024-bit integers
independently and uniformly at random until a prime is found.
What is the expected number of iterations before a prime is found?
What is the expected number of primes that are chosen before a
prime is encountered for the second time (i.e., the list of prime
numbers that are chosen contains two primes that are the same)?\\
\underline{Note~1}: Here, a 1024-bit integer is one whose binary representation
(with most significant bit equal to 1) has bitlength exactly 1024.\\
\underline{Note~2}: Let $\pi(x)$ denote the number of primes in $[2,x]$.
One version of the \emph{Prime Number Theorem} states that
$x/\log_e x < \pi(x) < 1.25x/\log_e x$ for $x \geq 114$.
You may use the approximation $\pi(x) \approx x/\log_e x$.\\
\underline{Remark}: This exercise should convince you that there are lots
of 1024-bit prime numbers. Also, if billions of people are selecting
1024-bit primes uniformly and independently at random, then the probability
that two people select the same prime is exceedingly small.
\end{problem}

\begin{loigiai}
\end{loigiai}
